% Added ARR experiments. Numeric sources are frozen; no model calls are needed.
% Stage 1: experiments/stage1_identity_safe_68_20260909_051438/formal/
%   final_summary.json, stage1_case_results_strict.csv, scc_equal_inference.json.
% Stage-1 manuscript values include the three author-confirmed corrections of 2026-10-13.
% These edits do not alter the underlying experiment files.
% Descriptive 80: stage1_legacy63_diagnostic/raw_input_manifest.json plus
%   the 17 contemporaneous cases above, replacing all four arms together.
% Stage 2: stage2_paired_replay_80case_identity_v1_final/
%   stage2_paired_core_table.csv, stage2_clustered_paired_inference.json.
% Cost: cost_prior_12case_gpt4o_gemini25pro_b60_v1/cost/strict_cost_table.csv.
% Prior: prior_fallback_slots_01_03_08_10/strict_prior_table.csv.
% Cost/Prior intervals: supplemental_pipeline_20260909/clustered_closeout/
%   paired_cluster_closeout.json. All paths above start at 2026NAACL/SA-MCGS/.
% Oracle definition: experiments/main_experiment/boosted_naive_baseline/
%   run_boosted_naive_round_robin.py, selected_metric_means/select_oracle_risk_top3.

\section{Additional Baseline, Ablation, and Cost Experiments}
\label{app:revision_experiments}

\subsection{Local Evidence Aggregation Protocol}
\label{app:revision_lea}

Local Evidence Aggregation (LEA) provides a structured baseline based on local graph retrieval and evidence aggregation, without implementing a full external GraphRAG system. For each SCC node, LEA constructs a four-node window from the node and its incoming and outgoing neighbors, using the fixed SCC node order. Missing positions are filled by alternating forward and backward offsets in that order. Each window is evaluated with the same local prompt builder, response schema, and parser as \samcgs{}, producing node scores, local risk nodes, conflict relations, critical pairs, and repair entries. Repeated evaluations may use cached results. The number of windows depends on graph size rather than a fixed rollout budget.

Let $\bar q_v$ denote the mean risk score over windows containing node $v$, and let $L_v$, $R_v$, and $E_v$ count its occurrences in local risk sets, repair entries, and conflict-edge endpoints, respectively. LEA ranks roots by $\bar q_v$ and assigns each node a subgraph membership score
\begin{equation}
 a_v=L_v+R_v+E_v+
 \max\{0,\bar q_v-\operatorname{median}_u\bar q_u\}.
\end{equation}
LEA returns the top $k=\min\{|S|,\max(4,\lceil0.45|S|\rceil)\}$ nodes, breaking ties by SCC order. If all membership scores are zero, it selects nodes by mean risk score. LEA uses no adaptive relation-memory selection, critical-pair revisiting, OC-guided core construction, or dynamic replacement. It therefore shares the local evidence evaluator with \samcgs{} but uses a simpler aggregation procedure. \revreplace{Computational cost is reported separately; token budgets are not matched.}{In this original five-model comparison, computational cost is reported separately and token budgets are not matched. Appendix~\ref{app:equal_budget} adds repeated-sampling baselines with per-case SA-referenced budgets.}

The main comparison comprises 80 injected variants from 20 source SCCs across Debian, SEC, BGB, and CUAD. GPT-4o, Gemini 2.5 Pro, Qwen2.5-72B, Llama 3.1-8B-Instruct, and Qwen3.5-9B each run One-shot, LEA, and \samcgs{} on the same inputs, yielding 400 model-cases and 1,200 method outputs. Inputs are matched across methods by source content, SCC membership, and injection identity. All Qwen3.5-9B runs use \texttt{enable\_thinking=false}. The separate contemporaneous cost experiment uses 12 source SCCs and 72 method jobs.

\paragraph{Oracle reference.}
The historical oracle-risk Top-3-from-10 reference uses ground-truth metrics to select among ten full-SCC attempts. For each three-attempt subset of valid outputs, it computes mean Root@3, Risk-any, and Risk-all. Selection lexicographically maximizes the minimum of these three metrics, their mean, and mean compression. With fewer than three valid attempts, it uses all valid attempts; with none, binary metrics are zero. Reported metrics average the selected attempts without pooling their nodes or subgraphs. This oracle selection requires ground truth, unlike deployable One-shot and LEA, and incurs the computation of ten attempts.

\begin{revadded}
\subsection{Repeated-Sampling Baselines with SA-Referenced Budgets}
\label{app:equal_budget}

\paragraph{Cases and budget.}
This comparison uses the same 80 injected variants from 20 source \scc{}s, with four variants per source: 40 CUAD, 20 SEC, 12 BGB, and 8 Debian cases. We evaluate GPT-4o, Qwen2.5-72B, Llama 3.1-8B-Instruct, Qwen3.5-9B, and Kimi-VL-A3B. Historical \samcgs{} outputs remain fixed. For model $m$ and case $i$, LEA and Full-SCC are assigned the same reference budget $B_{m,i}$ for candidate generation and selection. The reference uses recorded historical \samcgs{} token consumption where available, or calibration from existing usage records: same-model, same-domain SA cost measurements for GPT-4o, and historical SA call counts combined with same-case LEA tokens per call for Qwen2.5-72B.

\paragraph{Repeated generation and selection.}
Each LEA candidate reruns the complete four-node-window aggregation procedure in Appendix~\ref{app:revision_lea}; each Full-SCC candidate reanalyzes the entire component. Both retain their native root ranking and output core. The reference budget and selector context capacity determine the candidate allocation; the number of candidates is not fixed. With one candidate, it is returned directly. With multiple candidates, the same model acts once as a judge, reading the source records, graph relations, and anonymized candidates, including their rankings, cores, and up to 16 reported conflict pairs. It selects one existing candidate without merging or editing outputs. Candidate order is shuffled with a fixed seed, and selection never accesses injection labels or evaluation metrics.

\paragraph{Inference configuration.}
Candidate sampling uses temperature 0.4 and the judge uses temperature 0. Local-window, Full-SCC, and judge output limits are 2,048, 12,000, and 256 tokens, respectively, with request allowances adjusted to the reference budget. Qwen3.5-9B thinking is disabled, and the local vLLM context limit is 131,072 tokens. Relative to the original Full-SCC one-shot run, this protocol enables stochastic candidate sampling.

% Generated by build_table.py from archived five-model results.
\begin{table*}[t]
\centering
\renewcommand{\arraystretch}{1.02}
\small
\begin{tabular}{L{.23\textwidth}L{.13\textwidth}rrrr}
\toprule
\bfseries Model & \bfseries Method & \bfseries \rootthree{} & \bfseries \riskany{} & \bfseries \riskall{} & \bfseries Compression \\
\midrule
GPT-4o & LEA & \tblsecond{71.25} & \tblsecond{85.00} & \tblsecond{41.25} & \tblsecond{56.61} \\
 & Full-SCC & 22.50 & 53.75 & 23.75 & \tblbest{74.60} \\
\rowcolor{TableFocus}
 & \textbf{\samcgs{}} & \tblbest{86.25} & \tblbest{100.00} & \tblbest{75.00} & 51.26 \\
\midrule
Qwen2.5-72B & LEA & \tblsecond{55.00} & \tblbest{98.75} & \tblsecond{50.00} & \tblsecond{51.26} \\
 & Full-SCC & 23.75 & 85.00 & 46.25 & 45.32 \\
\rowcolor{TableFocus}
 & \textbf{\samcgs{}} & \tblbest{70.00} & \tblsecond{97.50} & \tblbest{75.00} & \tblbest{52.82} \\
\midrule
Llama 3.1-8B-Instruct & LEA & \tblsecond{50.42} & \tblsecond{93.75} & \tblsecond{34.58} & \tblsecond{51.22} \\
 & Full-SCC & 6.67 & 55.42 & 22.50 & 44.23 \\
\rowcolor{TableFocus}
 & \textbf{\samcgs{}} & \tblbest{61.25} & \tblbest{95.00} & \tblbest{48.75} & \tblbest{51.30} \\
\midrule
Qwen3.5-9B & LEA & \tblsecond{58.75} & \tblbest{93.33} & \tblsecond{62.08} & 51.26 \\
 & Full-SCC & 20.00 & 63.33 & 12.08 & \tblbest{71.76} \\
\rowcolor{TableFocus}
 & \textbf{\samcgs{}} & \tblbest{62.50} & \tblsecond{86.25} & \tblbest{63.75} & \tblsecond{67.25} \\
\midrule
Kimi-VL-A3B & LEA & \tblsecond{5.42} & \tblbest{89.58} & 17.50 & \tblsecond{51.26} \\
 & Full-SCC & \tblbest{7.92} & 59.58 & \tblbest{52.50} & 13.60 \\
\rowcolor{TableFocus}
 & \textbf{\samcgs{}} & 5.00 & \tblsecond{87.50} & \tblsecond{32.50} & \tblbest{51.85} \\
\bottomrule
\end{tabular}
\caption{Five-backbone comparison with SA-referenced baseline budgets on the same 80 variants per model (\%). \samcgs{} uses fixed historical outputs. Bold/underline mark the highest/second-highest value within each model, including ties. Shading identifies \samcgs{}; higher compression means a smaller core.}
\label{tab:equal_budget_five_models}
\end{table*}


\paragraph{Coverage and compactness.}
Table~\ref{tab:equal_budget_five_models} compares all five backbones. On GPT-4o, \samcgs{} reaches 86.25\% \rootthree{}, versus 71.25\% for LEA and 22.50\% for Full-SCC. On Qwen2.5-72B, it reaches 70.00\%, versus 55.00\% and 23.75\%, respectively, and its 52.82\% compression exceeds both baselines. For Qwen3.5-9B, Full-SCC achieves the highest compression (71.76\%) but only 12.08\% \riskall{}. Kimi shows the reverse balance: Full-SCC achieves the highest \riskall{} but only 13.60\% compression. These results show that the coverage--compression balance depends on the backbone.
\end{revadded}


\begin{table*}[!tp]
\centering
\renewcommand{\arraystretch}{1.10}

\small
\setlength{\tabcolsep}{6pt}
\begin{tabular}{lrrrrr}
\toprule
\bfseries Arm & \bfseries Root@3 $\uparrow$ & \bfseries Any $\uparrow$ & \bfseries All $\uparrow$ & \bfseries Comp. & \bfseries $|H|$ \\
\midrule
\multicolumn{6}{l}{\bfseries \stageoneupdate{80 variants / 20 source SCCs}} \\
Full \samcgs{} & \stageoneupdate{\tblbest{86.25}} & \tblbest{100.00} & \stageoneupdate{\tblbest{77.50}} & 51.26 & 8.10 \\
\relationmemorymark{No relation-first memory} & \relationmemorymark{\stageoneupdate{83.75}} & \relationmemorymark{\tblsecond{98.75}} & \relationmemorymark{\stageoneupdate{\tblsecond{73.50}}} & \relationmemorymark{51.26} & \relationmemorymark{8.10} \\
No critical-pair ledger/revisit & 83.75 & 88.75 & 56.25 & 71.29 & 4.30 \\
Random local windows & \stageoneupdate{\tblsecond{85.00}} & 96.25 & 71.25 & 51.26 & 8.10 \\
\bottomrule
\end{tabular}
\caption{\relationmemorymark{\stageoneupdate{Stage-1 ablations with GPT-4o on 80 variants from 20 source SCCs. Root@3, Any (Risk-any), All (Risk-all), and compression (Comp.) are percentages; $|H|$ is mean core size. Bold/underline: best/second-best endpoint metric, including ties; core size and compression are descriptive.}}}
\label{tab:revision_stage1}
\end{table*}

\subsection{\relationmemorymark{Stage-1 Component Ablations}}
\label{app:revision_stage1}

\relationmemorymark{\stageoneupdate{Table~\ref{tab:revision_stage1} compares Full, no relation-first memory, no critical-pair ledger/revisit, and random local windows on 80 variants from 20 source SCCs. Removing relation-first memory disabled relation-biased selection and relation probes; removing the critical-pair ledger disabled its protection and revisiting; random-window selection replaced the local-window selector. The table measures endpoint coverage in the returned core separately from root-ranking performance.}}

\relationmemorymark{\stageoneupdate{Full attained 77.50\% \riskall{}, compared with 73.50\% without relation-first memory, a difference of 4.00 percentage points at the same mean core size (8.10) and compression (51.26\%). Full also had higher \rootthree{} (86.25\% versus 83.75\%) and \riskany{} (100.00\% versus 98.75\%). These aggregate results favor retaining relation-first memory in this comparison.}}

\stageoneupdate{Removing critical-pair protection and revisiting reduced \riskall{} by 21.25 points, to 56.25\%, while reducing mean core size from 8.10 to 4.30. This comparison therefore combines changes in evidence selection and output size. Random local windows reduced \riskall{} by 6.25 points, to 71.25\%, at the same mean core size as Full.}

\relationmemorymark{\stageoneupdate{The Stage-1 results are descriptive comparisons across evaluation batches. We report aggregate differences without corresponding confidence intervals or significance tests.}}

\subsection{Stage-2 Paired Core Replay}
\label{app:revision_stage2}

We replayed frozen GPT-4o Full trajectories for 80 variants from 20 source SCCs without new model calls. Both policies received the same selector proposal at each trajectory prefix. The dynamic policy allowed node replacement; the locked policy admitted nodes up to the cap and retained them without eviction. Both shared OC signals, relation seeding, critical-pair protection, pair closure, thresholds, and the cap. All 80 pairs passed input-identity and final-Full parity checks.

\begin{table}[!htbp]
\centering
\renewcommand{\arraystretch}{1.10}

\small
\setlength{\tabcolsep}{4pt}
\begin{tabular}{lrrrrr}
\toprule
\bfseries Core policy & \bfseries Root@3 & \bfseries Any & \bfseries All & \bfseries Comp. & \bfseries $|H|$ \\
\midrule
Dynamic & \tblbest{86.25} & \tblbest{100.00} & \tblbest{75.00} & 51.26 & 8.10 \\
Locked & \tblbest{86.25} & \tblsecond{86.25} & \tblsecond{37.50} & 51.26 & 8.10 \\
\bottomrule
\end{tabular}
\caption{Stage-2 replay on shared search trajectories. Rates and compression are percentages. Root@3 is fixed by the shared trajectory and is unaffected by the core policy. \revadd{Bold/underline: best/second-best endpoint metric ($\uparrow$), including ties.}}
\label{tab:revision_stage2}
\end{table}

Dynamic replacement increased \riskall{} by 37.5 percentage points (95\% source-SCC paired interval $[12.5,60.0]$) and \riskany{} by 13.75 points ($[5.0,22.5]$), at equal mean core size and compression (Table~\ref{tab:revision_stage2}). The replay isolates endpoint retention given the observed evidence; search trajectories are fixed, and OC and closure are not ablated independently. \relationmemorymark{\stageoneupdate{The replay and Stage-1 evaluation (Table~\ref{tab:revision_stage1}) use different sets of Full outputs on the same 80 variants, yielding Full \riskall{} values of 75.00\% and 77.50\%, respectively.}}

% Generated by output/ablation_revision_20261011/build_appendix.py.
% Frozen values SHA256: 5a9ce882749c622d5ebe58a640b36c94d98d7d38117f413b2aad8229282a0287
\begin{ablationadded}
\subsection{Dynamic Core Replacement Across Backbones}
\label{app:new_core_ablation}

We extend paired core replay to Llama 3.1-8B-Instruct, Qwen3.5-9B, and Kimi-VL-A3B, with one Full run per model--case. Each model uses the same 80 injected variants: four templates applied to 20 source SCCs across CUAD, SEC, BGB, and Debian. Each run uses 60 rollouts, four-node windows, internal concurrency four, and temperature 0.4. These 240 trajectories form the ablation cohort, separate from the historical main-experiment results. Tables in this and the next three subsections report percentages averaged over 80 cases per model; the 20 source SCCs are the independent sampling units.

To isolate replacement, both policies receive the same ordered evidence proposals. Dynamic allows earlier admissions to be replaced; locked retains admitted nodes without eviction. For each case, locked capacity is fixed to the final Dynamic output size $k$, and both return exactly $k$ nodes. This paired replay isolates replacement at fixed evidence and retained-core size. Root@3 uses the shared trajectory ranking and is identical across selectors.

\begin{table}[!htbp]
\ablationcolor
\centering
\small
\renewcommand{\arraystretch}{1.10}
\setlength{\tabcolsep}{3pt}
\begin{tabular}{lrrrr}
\toprule
\bfseries Policy & \bfseries Root@3 & \bfseries Any & \bfseries All & \bfseries Comp. \\
\midrule
\multicolumn{5}{l}{\bfseries Llama 3.1-8B-Instruct} \\
Dynamic (Full) & \tblbest{62.50} & \tblbest{93.75} & \tblbest{62.50} & 51.26 \\
Locked & \tblbest{62.50} & \tblsecond{85.00} & \tblsecond{35.00} & 51.26 \\
\midrule
\multicolumn{5}{l}{\bfseries Qwen3.5-9B} \\
Dynamic (Full) & \tblbest{58.75} & \tblbest{88.75} & \tblbest{57.50} & 66.55 \\
Locked & \tblbest{58.75} & \tblsecond{70.00} & \tblsecond{33.75} & 66.55 \\
\midrule
\multicolumn{5}{l}{\bfseries Kimi-VL-A3B} \\
Dynamic (Full) & \tblbest{7.50} & \tblbest{93.75} & \tblbest{35.00} & 51.93 \\
Locked & \tblbest{7.50} & \tblsecond{80.00} & \tblsecond{30.00} & 51.93 \\
\bottomrule
\end{tabular}
\caption{Dynamic versus locked core replay at the default capacity ratio $r=0.45$. Any and All denote Risk-any and Risk-all; Comp. is $100(1-|H|/|S|)$. Evidence and actual output size are matched per case. Bold/underline mark the best/second-best endpoint rate within each model, including ties; compression is descriptive.}
\label{tab:new_core_replacement}
\end{table}

Dynamic replacement increases Risk-all by 27.50, 23.75, and 5.00 percentage points for Llama, Qwen, and Kimi, respectively (Table~\ref{tab:new_core_replacement}). Risk-any also improves on all three backbones at identical compression. The comparison supports revising the retained core as evidence accumulates, with the largest gains on Llama and Qwen.

\subsection{Core Capacity and Retention}
\label{app:new_core_capacity}

We replay the same Full evidence at four capacity ratios, $r\in\{0.25,0.35,0.45,0.60\}$, with cap $\min\{|S|,\max(4,\lceil r|S|\rceil)\}$. At each ratio, we recompute the production selector from the recorded evidence. Locked admission matches Dynamic's actual final $k$ case by case. The cap is an upper bound; observed compression is calculated from the actual retained-core size.

\begin{table}[!htbp]
\ablationcolor
\centering
\small
\renewcommand{\arraystretch}{1.10}
\setlength{\tabcolsep}{3pt}
\begin{tabular}{lrrrr}
\toprule
\bfseries $r$ & \bfseries Dynamic & \bfseries Locked & \bfseries $\Delta$ & \bfseries Comp. \\
\midrule
\multicolumn{5}{l}{\bfseries Llama 3.1-8B-Instruct} \\
0.25 & \tblbest{25.00} & 13.75 & $+11.25$ & 68.30 \\
0.35 & \tblbest{38.75} & 16.25 & $+22.50$ & 61.76 \\
0.45 & \tblbest{62.50} & 35.00 & $+27.50$ & 51.26 \\
0.60 & \tblbest{75.00} & 60.00 & $+15.00$ & 36.93 \\
\midrule
\multicolumn{5}{l}{\bfseries Qwen3.5-9B} \\
0.25 & \tblbest{47.50} & 22.50 & $+25.00$ & 73.49 \\
0.35 & \tblbest{48.75} & 26.25 & $+22.50$ & 70.31 \\
0.45 & \tblbest{57.50} & 33.75 & $+23.75$ & 66.55 \\
0.60 & \tblbest{57.50} & 38.75 & $+18.75$ & 64.21 \\
\midrule
\multicolumn{5}{l}{\bfseries Kimi-VL-A3B} \\
0.25 & \tblbest{17.50} & 7.50 & $+10.00$ & 68.34 \\
0.35 & \tblbest{20.00} & 13.75 & $+6.25$ & 61.91 \\
0.45 & \tblbest{35.00} & 30.00 & $+5.00$ & 51.93 \\
0.60 & \tblbest{55.00} & 47.50 & $+7.50$ & 38.67 \\
\bottomrule
\end{tabular}
\caption{Capacity sweep on the same 80 Full trajectories per model. Dynamic and Locked columns report Risk-all; $\Delta$ is their difference in percentage points. Comp. is the shared mean compression of the paired outputs. Bold marks the higher Risk-all within each pair.}
\label{tab:new_core_capacity}
\end{table}

Dynamic retains complete endpoints more often than locked admission at all 12 model--capacity settings (Table~\ref{tab:new_core_capacity}). Dynamic's Risk-all increases with capacity through $r=0.60$ for Llama and Kimi; Qwen reaches 57.50\% at both $r=0.45$ and $0.60$. The paired gains persist across the tested caps, while the absolute coverage and compression quantify the trade-off in output size.

\subsection{Selection from Shared Evidence}
\label{app:new_selector_control}

We compare the Full selector with LEA-score while holding the acquired evidence and final output size fixed. LEA-score applies the membership formula in Appendix~\ref{app:revision_lea} to the recorded Full evaluations and returns its top $k$ nodes, where $k$ is Full's actual output size for that case. Its signals are local-subgraph membership, repair-entry and conflict-endpoint counts, plus mean-risk surplus. The comparison isolates final selection on evidence acquired by SA-MCGS.

\begin{table}[!htbp]
\ablationcolor
\centering
\small
\renewcommand{\arraystretch}{1.10}
\setlength{\tabcolsep}{3pt}
\begin{tabular}{lrrrr}
\toprule
\bfseries Policy & \bfseries Root@3 & \bfseries Any & \bfseries All & \bfseries Comp. \\
\midrule
\multicolumn{5}{l}{\bfseries Llama 3.1-8B-Instruct} \\
Full & \tblbest{62.50} & \tblsecond{93.75} & \tblbest{62.50} & 51.26 \\
LEA-score & \tblbest{62.50} & \tblbest{96.25} & \tblsecond{57.50} & 51.26 \\
\midrule
\multicolumn{5}{l}{\bfseries Qwen3.5-9B} \\
Full & \tblbest{58.75} & \tblbest{88.75} & \tblbest{57.50} & 66.55 \\
LEA-score & \tblbest{58.75} & \tblsecond{86.25} & \tblsecond{52.50} & 66.55 \\
\midrule
\multicolumn{5}{l}{\bfseries Kimi-VL-A3B} \\
Full & \tblbest{7.50} & \tblbest{93.75} & \tblbest{35.00} & 51.93 \\
LEA-score & \tblbest{7.50} & \tblsecond{78.75} & \tblsecond{26.25} & 51.93 \\
\bottomrule
\end{tabular}
\caption{Full versus LEA-score at matched evidence and actual output size, using the default capacity. Endpoint-rate highlighting follows Table~\ref{tab:new_core_replacement}. Root@3 uses the shared trajectory ranking.}
\label{tab:new_selector_lea}
\end{table}

Full improves Risk-all over LEA-score by 5.00, 5.00, and 8.75 points on Llama, Qwen, and Kimi (Table~\ref{tab:new_selector_lea}).

\subsection{Evidence Accumulation Across Rollout Prefixes}
\label{app:new_prefix_analysis}

We reconstruct the native Full core after 15, 30, 45, and 60 completed rollout events from each stored 60-rollout trajectory. These four checkpoints measure evidence accumulation within each trajectory, with the selector, evidence thresholds, and nominal capacity ratio held fixed.

\begin{table}[!htbp]
\ablationcolor
\centering
\small
\renewcommand{\arraystretch}{1.10}
\setlength{\tabcolsep}{3pt}
\begin{tabular}{lrrrr}
\toprule
\bfseries Prefix & \bfseries Root@3 & \bfseries Any & \bfseries All & \bfseries Comp. \\
\midrule
\multicolumn{5}{l}{\bfseries Llama 3.1-8B-Instruct} \\
15 & 45.00 & 85.00 & 51.25 & 56.39 \\
30 & \tblbest{63.75} & \tblsecond{93.75} & \tblsecond{61.25} & 51.86 \\
45 & 57.50 & \tblbest{96.25} & 58.75 & 51.33 \\
60 & \tblsecond{62.50} & \tblsecond{93.75} & \tblbest{62.50} & 51.26 \\
\midrule
\multicolumn{5}{l}{\bfseries Qwen3.5-9B} \\
15 & 42.50 & 67.50 & 33.75 & 78.23 \\
30 & 55.00 & 76.25 & 45.00 & 73.19 \\
45 & \tblbest{60.00} & \tblsecond{86.25} & \tblsecond{51.25} & 69.02 \\
60 & \tblsecond{58.75} & \tblbest{88.75} & \tblbest{57.50} & 66.55 \\
\midrule
\multicolumn{5}{l}{\bfseries Kimi-VL-A3B} \\
15 & \tblsecond{5.00} & 73.75 & 13.75 & 63.14 \\
30 & \tblsecond{5.00} & 83.75 & 26.25 & 56.71 \\
45 & \tblbest{7.50} & \tblsecond{91.25} & \tblsecond{31.25} & 53.03 \\
60 & \tblbest{7.50} & \tblbest{93.75} & \tblbest{35.00} & 51.93 \\
\bottomrule
\end{tabular}
\caption{Native Full outputs reconstructed at four rollout prefixes. Each row averages the same 80 cases for its model. Root@3, Risk-any (Any), Risk-all (All), and compression (Comp.) are percentages. Bold/underline identify the highest/second-highest distinct endpoint rate across prefixes within each model; compression describes the evolving output size.}
\label{tab:new_native_prefixes}
\end{table}

From 15 to 60 events, native Risk-all rises from 51.25\% to 62.50\% on Llama, 33.75\% to 57.50\% on Qwen, and 13.75\% to 35.00\% on Kimi (Table~\ref{tab:new_native_prefixes}). The corresponding cores retain more nodes, quantifying the coverage--compression trade-off as search accumulates evidence. This trajectory analysis complements the paired replacement comparisons at fixed evidence and output size.

\end{ablationadded}


\begin{table*}[!tp]
\centering
\renewcommand{\arraystretch}{1.10}

\small
\setlength{\tabcolsep}{5pt}
\begin{tabular}{llrrrrrr}
\toprule
\bfseries Model & \bfseries Method & \bfseries Valid & \bfseries Invalid & \bfseries Root@3 & \bfseries Any & \bfseries All & \bfseries Comp. ($n$) \\
\midrule
GPT-4o & One-shot & 9 & 3 & 33.3 & 58.3 & 25.0 & 67.2 (9) \\
 & LEA & 12 & 0 & \tblsecond{83.3} & \tblsecond{91.7} & \tblsecond{41.7} & 51.0 (12) \\
\rowcolor{TableFocus}
 & \textbf{\samcgs{}} & 12 & 0 & \tblbest{91.7} & \tblbest{100.0} & \tblbest{83.3} & 51.0 (12) \\
\midrule
Gemini 2.5 Pro & One-shot & 12 & 0 & \tblbest{100.0} & \tblbest{100.0} & \tblbest{100.0} & 69.3 (12) \\
 & LEA & 12 & 0 & 75.0 & \tblbest{100.0} & 83.3 & 51.0 (12) \\
\rowcolor{TableFocus}
 & \textbf{\samcgs{}} & 12 & 0 & \tblsecond{91.7} & \tblbest{100.0} & \tblsecond{91.7} & 57.4 (12) \\
\bottomrule
\end{tabular}

\vspace{5pt}
\begin{tabular}{llrrrrrr}
\toprule
\bfseries Model & \bfseries Method & \bfseries Calls & \bfseries Input (K) & \bfseries Output (K) & \bfseries Total (K) & \multicolumn{2}{c}{\bfseries Time (s)} \\
\cmidrule(lr){7-8}
  &   &   &   &   &   & \bfseries Mean & \bfseries Median \\
\midrule
GPT-4o & One-shot & \tblbest{1.0} & \tblbest{17.0} & \tblbest{2.9} & \tblbest{19.9} & \tblbest{21.0} & \tblbest{22.1} \\
 & LEA & \tblsecond{17.8} & \tblsecond{36.3} & \tblsecond{22.8} & \tblsecond{59.1} & \tblsecond{189.2} & \tblsecond{138.0} \\
\rowcolor{TableFocus}
 & \textbf{\samcgs{}} & 59.6 & 112.3 & 68.4 & 180.7 & 570.9 & 574.7 \\
\midrule
Gemini 2.5 Pro & One-shot & \tblbest{1.0} & \tblbest{19.6} & \tblbest{7.5} & \tblbest{27.1} & \tblbest{85.9} & \tblbest{66.8} \\
 & LEA & \tblsecond{15.9} & \tblsecond{36.6} & \tblsecond{31.7} & \tblsecond{68.3} & \tblsecond{868.2} & \tblsecond{443.9} \\
\rowcolor{TableFocus}
 & \textbf{\samcgs{}} & 53.5 & 115.9 & 110.2 & 226.1 & 1650.2 & 1413.1 \\
\bottomrule
\end{tabular}
\caption{Cost, performance, and reliability on 12 source SCCs for each model--method pair. All binary rates use a denominator of 12, with invalid outputs counted as failures. Comp. is mean compression over valid outputs; parentheses give the denominator. Calls and tokens include recorded recovery usage and are divided by 12. Tokens are in thousands. Time is recorded runtime per case; parallel speedup was not measured. \revadd{Within each model, bold/underline mark best/second-best endpoint rates ($\uparrow$, upper panel) or costs ($\downarrow$, lower panel); compression and validity are descriptive. Shading identifies \samcgs{}.}}
\label{tab:revision_cost_full}
\end{table*}

\subsection{Contemporaneous Cost and Reliability}
\label{app:revision_cost}

The cost experiment used one injected variant from each of 12 distinct source SCCs. Debian/SEC/BGB/CUAD contributed 2/3/3/4 SCCs; each of the four templates contributed three variants; and small/medium/large SCCs numbered 6/3/3. GPT-4o and Gemini 2.5 Pro each ran One-shot, LEA, and \samcgs{} on these frozen inputs, yielding 72 method jobs. \samcgs{} used 60 rollouts, four-node windows, and internal concurrency one. One-shot generated a subgraph directly from the full SCC in a single call, without oracle selection. Table~\ref{tab:revision_cost_full} reports the results for each model--method pair.

Of the final outputs, 69 were valid and three were invalid. All invalid outputs came from GPT-4o One-shot and remained in the cost and performance denominators. One Gemini \samcgs{} job was interrupted and restarted, with 17 calls before interruption and 56 during recovery; its valid final output required more than one execution. Recorded usage includes both executions, although costs for unreturned in-flight requests are unknown. Runtime covers the recovery interval, with pre-interruption API time recorded separately.

The paired \samcgs{}--LEA \riskall{} differences were $+41.7$ points for GPT-4o (95\% interval $[8.3,75.0]$) and $+8.3$ for Gemini ($[-16.7,33.3]$). Gemini One-shot achieved the highest endpoint completeness in this sample. Averaged equally across models, \samcgs{} used 203.4K tokens per case, compared with 63.7K for LEA, or approximately $3.2\times$ as many. The performance gains from additional computation thus varied by model.

Internal rollout concurrency can reduce waiting time. A simple approximation to wall-clock time is
\begin{equation}
 T_{\mathrm{wall}}\approx T_{\mathrm{serial}}+
 \sum_g\left\lceil\frac{n_g}{p_g}\right\rceil t_g,
 \label{eq:revision_parallel_time}
\end{equation}
where each dependency-ordered group $g$ contains $n_g$ parallelizable calls, $p_g$ is effective parallelism, and $t_g$ is representative call latency. Available work, workers, and API limits constrain $p_g$. For fixed work, concurrency leaves calls, tokens, and fees unchanged. Equation~\ref{eq:revision_parallel_time} models execution time; no within-case speedup was measured, and concurrent execution of separate cases does not establish such a speedup.

\subsection{Relation-Prior Sensitivity}
\label{app:revision_prior}

We varied the relation-prior bonus multiplier (\texttt{relation\_prior\_weight}) from its default 0.85 to 0.68 and 1.02 ($\pm20\%$). The four source SCCs (cost slots 01/03/08/10) were selected and frozen before the sensitivity results were observed. Node, edge, and pair feature weights in Table~\ref{tab:prior_coefficients} remained fixed. The two models produced 16 new valid outputs; eight nominal outputs with matching inputs and configurations were reused from the cost experiment. All other settings were fixed, including rollout budget and internal concurrency one.

\begin{table}[!htbp]
\centering
\renewcommand{\arraystretch}{1.10}

\small
\setlength{\tabcolsep}{4pt}
\begin{tabular}{llrrr}
\toprule
\bfseries Model & \bfseries Metric & \bfseries 0.68 & \bfseries 0.85 & \bfseries 1.02 \\
\midrule
GPT-4o & Root@3 & \tblsecond{75} & \tblbest{100} & \tblsecond{75} \\
 & Risk-all & \tblbest{100} & \tblsecond{75} & \tblbest{100} \\
\midrule
Gemini Pro & Root@3 & \tblbest{100} & \tblsecond{75} & \tblbest{100} \\
 & Risk-all & \tblbest{75} & \tblbest{75} & \tblbest{75} \\
\bottomrule
\end{tabular}
\caption{Exploratory sensitivity to the relation-prior multiplier on four source SCCs per row. Root@3 and Risk-all are percentages; Risk-any is 100\% throughout. \revadd{Bold/underline: highest/second-highest distinct value across settings in each row, including ties.}}
\label{tab:revision_prior}
\end{table}

\begin{revadded}
\begin{table*}[!tp]
\centering
\small
\renewcommand{\arraystretch}{1.10}
\setlength{\tabcolsep}{3pt}
\begin{tabular}{L{0.21\textwidth}lrrr}
\toprule
\bfseries Model / scope & \bfseries Comparator & \bfseries Root@3 & \bfseries Risk-any & \bfseries Risk-all \\
\midrule
All five (fixed) & LEA & $+6.75\;[-0.25,14.00]$ & $-0.25\;[-4.50,4.00]$ & $+15.25\;[2.00,27.75]$ \\
 & One-shot & $+46.50\;[37.75,55.25]$ & $+35.25\;[24.00,46.75]$ & $+39.00\;[29.00,49.25]$ \\
\midrule
GPT-4o & LEA & $+7.50\;[-3.75,20.00]$ & $+3.75\;[0.00,10.00]$ & $+31.25\;[7.50,52.50]$ \\
 & One-shot & $+60.00\;[43.75,75.00]$ & $+45.00\;[27.50,62.50]$ & $+60.00\;[43.75,73.75]$ \\
\midrule
Gemini 2.5 Pro & LEA & $+0.00\;[-8.75,7.50]$ & $-2.50\;[-6.25,0.00]$ & $+0.00\;[-11.25,12.50]$ \\
 & One-shot & $-2.50\;[-16.25,11.25]$ & $+0.00\;[-5.00,5.00]$ & $-1.25\;[-20.00,16.25]$ \\
\midrule
Qwen2.5-72B & LEA & $+2.50\;[-11.25,16.25]$ & $+0.00\;[-6.25,5.00]$ & $+31.25\;[10.00,51.25]$ \\
 & One-shot & $+67.50\;[51.25,82.50]$ & $+51.25\;[31.25,71.25]$ & $+61.25\;[43.75,78.75]$ \\
\midrule
Llama 3.1-8B-Instruct & LEA & $+17.50\;[2.50,32.50]$ & $+2.50\;[-6.25,12.50]$ & $+13.75\;[-5.00,32.50]$ \\
 & One-shot & $+57.50\;[41.25,72.50]$ & $+47.50\;[30.00,65.00]$ & $+23.75\;[2.50,45.00]$ \\
\midrule
Qwen3.5-9B & LEA & $+6.25\;[-6.25,20.00]$ & $-5.00\;[-13.75,2.50]$ & $+0.00\;[-15.00,15.00]$ \\
 & One-shot & $+50.00\;[31.25,68.75]$ & $+32.50\;[18.75,47.50]$ & $+51.25\;[36.25,66.25]$ \\
\bottomrule
\end{tabular}
\caption{Main-experiment \samcgs{} minus comparator differences, in percentage points, with 95\% source-SCC bootstrap intervals in brackets.}
\label{tab:unified_paired_intervals}
\end{table*}
\end{revadded}

Endpoint completeness did not decrease on this subset, although root rankings changed (Table~\ref{tab:revision_prior}). This four-cluster check characterizes local sensitivity to one coefficient; it does not establish robustness to other priors or core thresholds. The sensitivity jobs ran concurrently across cases, precluding a direct runtime comparison with the serial nominal runs.



\subsection{Identity and Statistical Conventions}
\label{app:revision_provenance}

Frozen manifests distinguish physical source SCCs, injected variants, and model--method outputs. Inputs are matched by source content, SCC node membership, and injection identity, rather than traversal index alone. Experiment bundles retain raw outputs, result hashes, \stageoneupdate{available prompt/schema and configuration records}, parser versions, and call-level usage.

% Avoid a one-word continuation at the final column boundary.
\looseness=-1
\relationmemorymark{For the ablation and cost experiments, binary performance uses all planned outputs, counting invalid or undefined outcomes as failures. Compression is $1-|H|/|S|$ and is computed only for defined outputs, with the denominator reported. Risk-any and Risk-all measure intersection and full inclusion of the injected endpoint set in $H$, respectively, rather than node-level precision. \stageoneupdate{Where intervals are reported, they use 20,000 percentile-bootstrap replicates with seed 8229, resampling physical source SCCs jointly across methods.} Source-equal estimates average variants within each SCC before averaging SCCs; variant-level averages are labeled explicitly. Models and variants from the same source do not constitute independent clusters; statistical comparisons must account for such dependence \citep{dror2018hitchhiker}. The intervals characterize the selected source set, not natural-risk prevalence.}

\begin{revadded}
% Keep the final subsection with its full-width interval table in the clean view.
\ifrevisionmarks\else\newpage\fi
\subsection{Main-Experiment Confidence Intervals}
\label{app:main_confidence_intervals}

For the main experiment, we resampled 20 source SCCs with replacement as paired clusters, keeping each source's four variants, five fixed models, and methods together. Paired differences are averaged within each source over variants and, for the overall estimate, models; source means are then averaged equally. Table~\ref{tab:unified_paired_intervals} reports overall and model-specific 95\% percentile intervals from 20,000 bootstrap replicates (seed 8229). These intervals reflect variation across the selected sources rather than repeated-run variability.


\end{revadded}

\begin{revadded}
\subsection{Computational Resources and Implementation Details}
\label{app:compute_implementation}

\paragraph{Computational resources and budget.}
Local inference was conducted on a cluster equipped with eight NVIDIA H800 GPUs, with GPU allocation depending on the model. GPT-4o and Gemini 2.5 Pro were accessed through APIs. Across evaluation and development experiments, we estimate aggregate inference usage to be approximately one billion input and output tokens, including recorded retries. Available usage logs account for approximately 0.95 billion tokens; missing historical usage is estimated from comparable cost profiles. Offline analyses reuse existing trajectories without additional model inference.

\paragraph{Software implementation and configuration.}
The implementation uses Python and NetworkX for graph processing, with QuantLaw for BGB reference extraction. Recorded local deployments use vLLM 0.20.0, PyTorch 2.11.0, and Transformers 5.6.2. The default \samcgs{} configuration uses 60 rollout slots and four-record windows. Evaluation metrics and source-level bootstrap confidence intervals are computed using custom scripts; experiment-specific settings are provided in the corresponding appendix subsections.
\end{revadded}
